\begin{abstract}

    Die zunehmende Urbanisierung und Motorisierung führen zu einer wachsenden Belastung städtischer Verkehrsinfrastrukturen, wobei die Parkplatzsuche eine zentrale Rolle spielt. \acp{sps} auf Basis von \ac{iot}-Technologien und Cloud-Architekturen gewinnen in diesem Kontext zunehmend an Bedeutung. Trotz ihres Potenzials zur Reduktion von Suchverkehr, Emissionen und Lärm stehen Parkhausbetreiber vor der Herausforderung, die wirtschaftliche Rentabilität solcher Systeme verlässlich zu bewerten -- insbesondere die laufenden Kosten für Sensorik und Cloud-Dienste im Verhältnis zum erwarteten Mehrwert.
    
    Um diese Lücke zu adressieren, wird in dieser Arbeit der Ansatz \enquote{Drive \& Decide} entwickelt und evaluiert. Dazu wird ein prototypisches, \ac{iot}-basiertes \ac{sps} mit serverless \ac{aws}-Architektur umgesetzt, das ein physisches Parkmodell mit einer virtuellen Lastsimulation verbindet. Darauf aufbauend wird ein Kosten-Nutzen-Modell entwickelt, das sowohl die einmaligen \ac{capex} als auch die nutzungsabhängigen \ac{opex} der Cloud-Infrastruktur systematisch erfasst und über \ac{tco}, \ac{roi}, Amortisationsdauer und Break-Even-Punkt bewertet.
    
    Die Evaluierung bestätigt die technische Machbarkeit des Prototyps und zeigt, dass der durchgehend ereignisgetriebene, serverless Aufbau nur geringe laufende Cloud-Kosten verursacht. Als belastbarste Kennzahl erweist sich der Break-Even-Punkt: Bereits eine Auslastungssteigerung von rund 0{,}03 bis 0{,}04 Prozentpunkten genügt, damit sich die Investition über fünf Jahre trägt -- deutlich unterhalb der in der Literatur berichteten Effekte. Damit liefert der Ansatz eine datengetriebene und nachvollziehbare Entscheidungsgrundlage und zeigt, dass ein solches \ac{sps} unter den getroffenen Annahmen wirtschaftlich tragfähig ist.
    
\end{abstract}